\section{Discussion}
\label{sec:discussion}

% TODO: full discussion. The published paper will discuss:
%   - The platform vs. the empirical question (separation of concerns)
%   - Domain shift from PSG to Muse S
%   - The 4-class sleep staging limit
%   - LLM as a moving target
%   - Privacy posture (on-device, no telemetry)
%   - Open-source contribution vs. closed-system research

The NeuralCompose platform is useful independently of whether the D8
hypotheses are supported. If the active condition produces a significant
increase in novelty, the platform is a validated research instrument. If
it does not, the platform is a documented open-source failure mode with
well-instrumented data and a pre-registered design --- itself a
contribution, since pre-registered negative results are publishable when
they are well-instrumented.

The largest expected source of model error is the channel mapping from
public PSG datasets (central and occipital derivations) to Muse S (4
unipolar frontal channels). We expect a model trained on Sleep-EDFx to
underperform on per-user Muse S data, and we plan per-user fine-tuning
as the path to acceptable accuracy. The validation toolkit provides the
debugging surface for this iteration.

The 4-class output is the honest upper bound for sleep staging from Muse
S alone. We do not claim 5-class AASM scoring. The \texttt{Uncertain\_REM}
flag is a weak REM proxy and is labeled accordingly.

The local LLM is a moving target. The platform specifies the protocol
(\texttt{DreamAnalysisPredicting}) and the prompt templates; the
underlying model is swappable. We document the model version per session
to make the analyses reproducible.

The privacy posture is structural: no network calls at runtime, no
telemetry, no cloud sync. All session data, EEG, dream reports, and
analyses stay on the user's machine. A single action deletes everything.
This is not a feature; it is a constraint that the codebase enforces.
